% title: Finite Sets with a Divisibility Closure Property
% short-title: Finite Divisibility Sets
% abstract: We formalize in Lean 4 the classification of finite nonempty sets of positive integers satisfying the condition that for every pair of elements there is a third element giving the divisibility relation used in the problem. The proof is structural: parity forces a common power of two factor, and in the odd case the map x goes to (M minus x) over 2 on the non-maximal elements collapses the set to the two possibilities {d} and {d,3d}.
% subjclass: 11A05
% keywords: divisibility, finite sets, olympiad problem, Lean 4
% author: Mario Hernández M.

\subsection{Finite Sets with a Divisibility Closure Property}

We formalize the hypothesis as
\lean{Biblioteca.Demonstrations.DivisibilitySet} and prove the classification
theorem
\lean{Biblioteca.Demonstrations.divisibility_set_classification}.

\begin{theorem}
Let \(S\) be a finite nonempty set of positive integers such that for every
\(a,b \in S\) there exists \(c \in S\) with
\[
  a \mid b + 2c.
\]
Then \(S\) is either
\[
  \{d\}
  \qquad\text{or}\qquad
  \{d,3d\}
\]
for some positive integer \(d\). Conversely, every set of one of these two
forms satisfies the divisibility condition.
\end{theorem}

\begin{proof}
The converse is immediate. If \(S=\{d\}\), then choosing \(c=d\) gives
\(d \mid d+2d\). If \(S=\{d,3d\}\), then for \(a=d\) the divisibility is
obvious, while for \(a=3d\) one chooses \(c=b\), so that
\[
  3d \mid b+2b = 3b.
\]

For the classification, first observe that if \(S\) contains an even element
\(a\), then every element of \(S\) is even: indeed \(a\mid b+2c\) forces
\(b+2c\) to be even, hence \(b\) is even. Therefore either all elements are
odd or all are even. In the even case we divide the whole set by \(2\); the
same divisibility relation is preserved, so by iterating this reduction we may
assume that all elements of \(S\) are odd.

Now let \(M=\max S\). If \(S=\{M\}\), we are done. Otherwise consider any
\(x \in S\) with \(x<M\). The hypothesis for \(a=M\) gives some \(c \in S\)
such that \(M \mid x+2c\). Since \(x\), \(c\), and \(M\) are odd, the number
\(x+2c\) is odd, so it cannot equal \(2M\). Also \(c \neq M\), because that
would imply \(M \mid x\) with \(0 < x < M\). Hence \(c<M\), and therefore
\[
  0 < x+2c < 3M.
\]
The only positive odd multiple of \(M\) in that range is \(M\) itself, so
\[
  x+2c = M.
\]
Thus the map
\[
  f(x) = \frac{M-x}{2}
\]
sends \(S \setminus \{M\}\) to itself.

Let \(m\) and \(v\) be respectively the minimum and maximum of
\(S \setminus \{M\}\). The map \(f\) is injective on this finite set, hence it
permutes it. Since \(f\) is strictly decreasing, its image of the minimum is
the maximum, and its image of the maximum is the minimum:
\[
  f(m)=v,\qquad f(v)=m.
\]
Equivalently,
\[
  2v = M-m,\qquad 2m = M-v.
\]
Subtracting these equations gives \(v=m\), and then \(M=3m\). So every
non-maximal element is equal to \(m\), and therefore
\[
  S = \{m,3m\}.
\]
Undoing the powers of \(2\) removed at the beginning yields the final
classification
\[
  S=\{d\}\quad\text{or}\quad S=\{d,3d\}.
\]
\end{proof}
